我们以两种形式提供预先构建好的模拟器：Python 代码和容器文件。
容器文件由 Python 代码生成；如果学生要运行 Python 代码，则需要从 GitHub
安装 SEED Emulator 源代码。容器文件可以直接使用，不需要安装模拟器源代码。
如果教师希望定制模拟器，可以修改 Python 代码，生成自己的容器文件，然后把
这些文件提供给学生，用来替换实验设置文件中包含的文件。具体说明见
\texttt{README.md} 文件。


\paragraph{下载模拟器文件。}
请从本实验网页下载 \texttt{Labsetup.zip} 文件并解压。容器文件夹中的文件是
实际使用的模拟文件；它们由 Python 代码生成。对于大多数实验，容器文件夹名为
\texttt{output/}；如果某个实验包含多个模拟器，则会使用不同的文件夹名。实际
名称会在实验任务中给出。


\paragraph{启动模拟器。}
我们将直接使用容器文件夹中的文件。进入该文件夹，运行 Docker 命令来构建并
启动容器。我们建议在提供的 SEED Ubuntu 20.04 VM 中运行模拟器；只要安装了
Docker 软件，在普通 Ubuntu 20.04 操作系统中运行通常也没有问题。读者可以从
\href{https://github.com/seed-labs/seed-labs/blob/master/manuals/docker/SEEDManual-Container.md}
{\underline{这个链接}}查看 Docker 使用手册。如果这是你第一次使用容器搭建
SEED 实验环境，请务必阅读用户手册。


下面列出了一些常用的 Docker 和 Compose 命令。由于这些命令会被频繁使用，
我们已经在提供的 SEED Ubuntu 20.04 VM 的 \texttt{.bashrc} 文件中为它们
创建了别名。

\begin{lstlisting}
$ docker-compose build  # Build the container images
$ docker-compose up     # Start the containers
$ docker-compose down   # Shut down the containers


// Aliases for the Compose commands above
$ dcbuild       # Alias for: docker-compose build
$ dcup          # Alias for: docker-compose up
$ dcdown        # Alias for: docker-compose down
\end{lstlisting}


所有容器都会在后台运行。要在某个容器中运行命令，我们通常需要进入该容器的
shell。首先使用 \texttt{"docker ps"} 命令找出容器 ID，然后使用
\texttt{"docker exec"} 在该容器中启动一个 shell。我们也已经在
\texttt{.bashrc} 文件中为这些命令创建了别名。

\begin{lstlisting}
$ dockps        // Alias for: docker ps --format "{{.ID}}  {{.Names}}" 
$ docksh <id>   // Alias for: docker exec -it <id> /bin/bash

// The following example shows how to get a shell inside hostC
$ dockps
b1004832e275  hostA-10.9.0.5
0af4ea7a3e2e  hostB-10.9.0.6
9652715c8e0a  hostC-10.9.0.7

$ docksh 96
root@9652715c8e0a:/#  

// Note: If a docker command requires a container ID, you do not need to 
//       type the entire ID string. Typing the first few characters will 
//       be sufficient, as long as they are unique among all the containers. 
\end{lstlisting}


如果在搭建实验环境时遇到问题，请阅读手册中的 ``Common Problems'' 部分，
寻找可能的解决办法。


\begin{comment}
\paragraph{设置终端标题。}
我们可能需要进入多个容器，并打开多个终端标签页，在它们之间来回切换。这样
很容易混淆哪个标签页对应哪个容器。为了解决这个问题，进入容器后可以使用
下面的命令设置终端标题；命令会把标题设置为 \texttt{"New Title"}。

\begin{lstlisting}
# set_title New Title
# st New Title       (*@\pointleft{st} is an alias of set\_title@*)
\end{lstlisting}
\end{comment}

